\subsection{An action on the entire mixed Fermat homology tower}

\begin{statement}{D.39}{Proved}{An integral all-prime homology lift retaining the degree-eight and CM data}
Use the marked groups $B,F,K,L,L_4,H,H_4$ of D.34--D.38,
and write closures inside $\widehat B$. Set
\[
 \begin{split}
 \mathcal H&=\widehat B/\overline{[L,L]}
       =\varprojlim_m B/[L,L]L^m,\\
 \mathcal T&=\widehat B/
       \bigl(\overline{[L,L]}\cap\overline{D_9K}\bigr),
       \qquad T_3=B/D_9K.
 \end{split}                                                   \tag{D.39.1}
\]
The maximal common marked quotient of $\mathcal H$ and $T_3$ is
\[
 C_8=B/\bigl([L_4,L_4]D_9K\bigr),\qquad
 \mathcal T\simeq\mathcal H\mathbin{\times_{C_8}}T_3.
                                                               \tag{D.39.2}
\]
Its coefficient kernel and order are
\[
 \ker(C_8\to H_4)\simeq C_8^7\oplus C_4^{62}\oplus C_2^{265},
 \qquad |C_8|=2304\,2^{410}=18\,2^{417}.                       \tag{D.39.3}
\]
Here $C_n$ in an abelian group type denotes a cyclic group of
order $n$. The canonical map from this common quotient to D.38's
common quotient has kernel $C_2^7$.

There is an actual continuous automorphism of $\mathcal T$,
with multiplier exactly one, satisfying the peripheral and two
hexagon equations in this quotient. It gives one compatible generating family of
finite GT0 shadows throughout this homology tower. Its normalized
$G_6$ and $T_3$ actions are exactly those of D.38. Its normalized
action on every odd-primary homology component
$\varprojlim_n B/[L,L]L^{p^n}$ is the identity.
In particular, it replaces D.38's $J_3$ and $J_{12}$ actions,
while retaining the fixed-fibre pair $T_4=I,\ W_3=[i]$.

The construction retains an explicitly checked seed on
\[
 \mathcal J_{48}=J_{48}\mathbin{\times_{C_8}}T_3,
 \qquad J_{48}=B/[L,L]L^{48},                               \tag{D.39.4}
\]
with
\[
 \begin{gathered}
 |\mathcal J_{48}|=2^{15771}3^{1156},\qquad
 |\ker(\mathcal J_{48}\to J_{48})|=2^{11151},\\
 |\ker(\mathcal J_{48}\to T_3)|=2^{4202}3^{1154}.
 \end{gathered}                                               \tag{D.39.5}
\]
Its cusp has order $1152$ and its finite-shadow $N_{\rm ord}$
is $576$. The recorded seed has $\beta^{\rm ab}=(1482,1044)$
and gauge $k=-438$. Its normalized degree-eight action is
\[
 \operatorname{Int}(c_0^{146})\alpha_{33}
   =\operatorname{Int}(c_0^{258})\alpha_{33},                 \tag{D.39.6}
\]
the second expression being D.38's action. The equality follows
from the exact order $16$ of $c_0=t^6$ in $T_3$.
Through the marked compact-cover maps of D.36 and D.38, its
actual matrices are
\[
 \rho_{C_3,16}=\begin{pmatrix}3&8\\8&11\end{pmatrix},
 \qquad \rho_{C_4,3}=I.                                     \tag{D.39.7}
\]
Thus every finite refinement in this family retaining
$\mathcal J_{48}$ has the same nonarithmetical CM data.

The exact remaining kernel of each source map
\[
 \widehat B\twoheadrightarrow\mathcal T,\qquad
 \widehat F\twoheadrightarrow\ker(\mathcal T\to S_3)
\]
is $\overline{[L,L]}\cap\overline{D_9K}$, with the open
embedding of $\widehat F$ understood. It contains the nonzero
word $[a_1^{12},a_2^{12}]^2$ of D.38. The canonical projection
from $\mathcal T$ to D.38's underlying finite joint group is
strict: $x^{144}$ is killed there and survives in $\mathcal T$.
The present action induces a replacement of its $J_{12}$
component, as specified above. D.45 proves that this recorded action has no lift to an
actual full profinite GT0 element. The prescribed CM pair
remains a target for a different action.

This addresses Part D(e), section 51.16, folio 734 verso,
scan 174, and D-after-LM2(1), section 51.15, folios 725--727,
scans 156--160. It establishes compatibility across every
prime and every finite homology level, together with the
retained degree-eight action.
\end{statement}
\begin{proof}
We first specify the integral source and its prime components.
The finite-index subgroup $L$ is free of rank $1153$.
Its induced profinite topology is its full topology: for a
finite-index subgroup of $L$, its normal core in $B$ is again
finite-index. Thus its closure is $\widehat L$, and
\[
 \overline L/\overline{[L,L]}\simeq\widehat{\mathbb Z}^{1153}.
                                                               \tag{D.39.8}
\]
The graph of D.37 has $6912$ edges and $5760$ vertices. Its
integer cycle matrix has the identity on its $1153$ chosen
non-tree edges. Consequently its path model identifies
$B/[L,L]$ with the subgroup of $\mathbb Z[H]\rtimes H$
whose elements have the specified path boundary. Reducing
the coefficients modulo $m$, or completing them in
$\mathbb Z_p$ or $\widehat{\mathbb Z}$, gives the corresponding
groups without losing any cycle coordinate.

Put $\mathcal H_p=\varprojlim_nJ_{p^n}$. There is a canonical
isomorphism
\[
 \mathcal H\longrightarrow\prod_{p,H}\mathcal H_p,          \tag{D.39.9}
\]
where the product consists of tuples with the same image in
$H$. To prove surjectivity, choose one word lifting their
common base element and subtract its path in every component;
the remaining cycle vectors form an element of
$\prod_p\mathbb Z_p^{1153}=\widehat{\mathbb Z}^{1153}$.
The same cycle-coordinate description proves injectivity.
This also proves the inverse-limit assertion in (D.39.1).

\paragraph{The stabilized common quotient.}
Work first in
$E_8=B/[L_4,L_4]L_4^8$, with coefficient subgroup
$V_8=(\mathbb Z/8)^{385}$, and let $G$ be the image of $K$.
The quotient $A_0=K/L_4$ has order $128$, class two and
exponent four; its derived group has order two. These facts
are checked on all its elements using the faithful
$H_4$ graph model, and also follow from D.38's embedding
in $O_2(Q)\times(\mathbb Z/4)^2$ together with the displayed
generators. Hence $\gamma_3G$ lies in the abelian subgroup
$V_8$. With $I$ the augmentation ideal of the actual four
$K$-generator actions and $U=\gamma_3G$, one has
$\gamma_jG=I^{j-3}U$ for $j\ge3$.

The full Zassenhaus product in this exponent-eight module is
\[
 W_8:=\operatorname{im}(D_9K\to E_8)
   =G^{16}+(\gamma_2G)^8+4\gamma_3G
        +2\gamma_5G+\gamma_9G.                            \tag{D.39.10}
\]
Each displayed summand lies in $V_8$, so the sum is an
abelian-module sum. To exhaust the product, the terms with
$i=1,2$ begin at powers $16,8$, respectively; those with
$i=3,4$ are contained in $4\gamma_3G$; those with
$5\le i\le8$ are contained in $2\gamma_5G$; and all terms
with $i\ge9$ lie in $\gamma_9G$. Further powers are already
in these summands. Surjectivity of $K\to G$ lifts every
commutator and power appearing in this description.

The two power summands in (D.39.10) are computed on their
entire groups. Choose actual lifts $g_h$ of all $128$
elements $h$ of $A_0$, and denote their coefficient actions
by $M_h$. Every element of $G$ is $v g_h$, with $v\in V_8$,
and
\[
 (v g_h)^{16}=(1+M_h+\cdots+M_h^{15})v+g_h^{16}.
\]
For $h$ of order one or two the norm vanishes modulo eight.
For $h$ of order four it is
$4(1+M_h+M_h^2+M_h^3)$. The program includes each actual
$g_h^{16}$ and every column of this norm on the full
coefficient module. Their span is $C_2^{327}$.

For the derived subgroup, select an actual basic commutator
$c$ projecting to the nontrivial element of $A_0'$. Its
intersection with $V_8$ is generated by the following
explicit vectors: $c^2$; the conjugates of the six basic
$K$ commutators by all $g_h$, subtracting $c$ when their
base image is nontrivial; and $(1-M_c)V_8$. These vectors
belong to $\gamma_2G$. Conversely their span, together
with the coset represented by $c$, is a normal subgroup
containing every basic commutator. Normality under the
four generators and $M_c$ is checked, and conjugation
by $V_8$ is included by $(1-M_c)V_8$. This proves equality
with $\gamma_2G$, including all conjugator freedom. Its
coefficient intersection has type $C_8^{381}$. The eighth
powers in the nontrivial coset are exactly
\[
 (v c)^8=4(1+M_c)v+c^8,
 \qquad v\in\gamma_2G\cap V_8;
\]
the eighth powers in the identity coset vanish. Their span
is $C_2^{189}$.

The remaining summands are obtained by exact module
elimination over $\mathbb Z/8$. The successive group types
are as follows; the entries give the numbers of cyclic
factors of the indicated orders.
\[
 \begin{array}{c|rrr}
 j& C_8&C_4&C_2\\\hline
 3&375&1&0\\
 4&355&11&10\\
 5&316&38&18\\
 6&252&85&18\\
 7&178&134&28\\
 8&106&146&84\\
 9&51&127&134
 \end{array}                                                   \tag{D.39.11}
\]
The normal closure of the twenty-four basic triple
commutators gives the first row. Repeated multiplication
by $I$ gives the subsequent rows. Every translate by a
marked $K$ generator is checked. Summing all five terms
in (D.39.10) gives
\[
 W_8\simeq C_8^{51}\oplus C_4^{265}\oplus C_2^{62}.         \tag{D.39.12}
\]
Its full $s,r$ stability is separately verified. Reduction
modulo four gives exactly D.38's recorded $D_9$ image.
Quotienting $V_8$ by (D.39.12) proves (D.39.3), and the
kernel of reduction to the earlier common quotient is
the quotient of $4V_8$ by its intersection with $W_8$.
It has order $2^7$ and exponent two.

These finite computations determine the common quotient
at every higher precision. Indeed
$L_4\subset N\subset D_2K$, whence
\[
 L_4^8\subset D_{16}K\subset D_9K.                        \tag{D.39.13}
\]
In $\mathcal E_2=\varprojlim_n B/[L_4,L_4]L_4^{2^n}$,
the image of $D_9K$ therefore contains
$8\mathbb Z_2^{385}$ and is exactly the inverse image
of $W_8$ modulo eight. D.37's relative-prime argument
also applies integrally: the maximal relative pro-two
quotient of $\mathcal H$ over $Q$ is $\mathcal E_2$.
Explicitly, the complementary group $L_4/L\simeq C_3$
has vanishing positive-degree homology over $\mathbb Z_2$,
and its coinvariants on $L^{\rm ab}\otimes\mathbb Z_2$
are $L_4^{\rm ab}\otimes\mathbb Z_2$. Killing its normal
closure and the odd coefficient factors gives this map;
every quotient over $Q$ with pro-two kernel kills them
and factors through the same map. Since $T_3\to Q$ has
two-group kernel, every common marked quotient factors
through $\mathcal E_2$. Dividing it by the just-computed
image of $D_9K$ proves (D.39.2) in both directions.
The quotient by the intersection of the two closed
normal kernels then gives its asserted fibre product.

\paragraph{The exact integral two-adic correction.}
We retain the actual conjugators of D.33 in $T_3$.
They are reconstructed from their recorded coefficients;
both SHA256 values and their $12\,921\,196$-letter cusp
reconstruction agree with the original certificate.
Their images in $E_8$ determine their exact common images
in $C_8$. Their cusp defect belongs to $W_8$, which is
checked independently of the next solver.

Let $A_H,\beta_H$ be D.37's chosen conjugator images in
$H$, let $P_A,P_\beta$ be their canonical integral
paths, and put $s_*=\beta_Hs\beta_H^{-1}$ and
$r_*=A_HrA_H^{-1}$. Thus $s_*r_*=t$. Left corrections
by cycle vectors $a,b\in\mathbb Z_2^{1153}$ give the
complete cusp equation
\[
 (1-s_*)b+(s_*-t)a=-E,                                   \tag{D.39.14}
\]
where $E$ is the integral defect of the two canonical
paths. All terms are written in the identity-minor
cycle basis. This equation is exact because the
coefficient subgroup is abelian; it includes the full
chosen paths and their base action.

Let $S_4$ and $S_Q$ be the integer cycle maps induced by
$H\to H_4$ and $H\to Q$. Let $\mathcal W\subset
\mathbb Z_2^{385}$ be the inverse image of $W_8$.
Its basis is the $378$ recorded rows of (D.39.12),
together with $8e_j$ on the seven nonpivot columns.
After the recorded column permutation, its basis
matrix is triangular with diagonal
\[
 1^{51},\quad 2^{265},\quad4^{62},\quad8^7.
                                                               \tag{D.39.15}
\]
It has index $2^{410}$ and contains $8\mathbb Z_2^{385}$.
Write $Z$ for its column matrix. In addition to
(D.39.14), impose
\[
 \begin{array}{ll}
 S_4b-Zu_\beta=r_\beta,&S_4a-Zu_A=r_A,\\
 S_Qb-2v_\beta=r_{\beta,2},&S_Qa-2v_A=r_{A,2}.
 \end{array}                                                   \tag{D.39.16}
\]
The first row prescribes the actual D.33 conjugators in
$C_8$. The second row prescribes exactly the retained
$G_2$ conjugators. Each right-hand side is the difference
between the target path and the canonical path under
the indicated map, in its recorded integer residue
representative. The variables $u_A,u_\beta$ have $385$
coordinates and $v_A,v_\beta$ have $97$ coordinates.
Thus (D.39.14)--(D.39.16) form a signed integer matrix
$M$ of size $2117\times3270$.

The exact elimination modulo $256$ gives full row rank,
with invariant factors
\[
 1^{1588},\quad2^{457},\quad4^{64},\quad8^7,\quad16.
                                                               \tag{D.39.17}
\]
These determine the integral $\mathbb Z_2$ invariant
factors. To see why the finite calculation suffices,
lift its elementary row and column operations to
$\mathbb Z_2$. The transformed matrix is
$[D\ 0]+256R$, with $D$ the diagonal in (D.39.17).
Factoring $D$ on the left leaves a square leading block
congruent to the identity modulo $16$, hence invertible.
The remaining columns are also divisible by the
corresponding diagonal entries. Thus the exact image
has the displayed invariant factors. In particular
\[
 16\mathbb Z_2^{2117}\subset M\mathbb Z_2^{3270},
 \qquad \operatorname{rank}_{\mathbb Z_2}\ker M=1153.
                                                               \tag{D.39.18}
\]

The recorded solution $z$ modulo $256$ has a full-rank
$1153\times1153$ kernel map modulo two. It lifts to an
exact solution retaining $z$ modulo $16$: its integral
residual is divisible by $256$, so dividing that residual
by $16$ places it in the image in (D.39.18). This proves
existence without assuming that an arbitrary mod-two
solution lifts. The certificate also constructs the lift
recursively. From a solution modulo $2^{k+4}$, solve
\[
 M\delta=\frac{b-Mz}{2^k}\pmod{256},\qquad
 z_{\rm new}=z+2^k\delta\pmod{2^{k+8}},                    \tag{D.39.19}
\]
where $b$ denotes the whole signed right-hand side.
The new right-hand side is divisible by $16$, so
(D.39.18) makes this finite equation solvable. Beginning
with $k=4$ and increasing $k$ by four gives compatible
coefficients with a two-adic limit. Two exact steps,
retaining moduli $16$ and $256$ and solving at moduli
$4096$ and $65536$, are checked using the actual signed
integer matrix. The multiplier remains exactly one.

The resulting homomorphism of $\mathcal H_2$ induces
the prescribed base automorphism and a matrix in
$\operatorname{GL}_{1153}(\mathbb Z_2)$ on its coefficient
kernel. The mod-two invertibility just verified proves
this integral invertibility. It is therefore an actual
automorphism. Equations (D.39.16) retain exactly the
common conjugators needed to pair it with $T_3$.

\paragraph{One formula for every odd prime.}
The selected base action on $H$ is precisely
$\operatorname{Int}(x^6)$, checked on all $6912$ elements;
moreover $A_H=x^6$. Put $h=x^{-6}\beta_H$. It centralizes
$s$ in $H$. Choose its canonical word $c$ and path $P_c$.
The vector $\delta=P_c-sP_c$ is an integral cycle and
satisfies $s\delta=-\delta$. At every odd prime, replace
the path of $c$ by
\[
 P_{c,p}=P_c-\tfrac12\delta
          =\tfrac12(P_c+sP_c).                            \tag{D.39.20}
\]
This is an element of $\mathcal H_p$ with the same base
image, and it centralizes $s$ exactly. Indeed its base
commutes with $s$, its path is $s$-fixed, and $P_s=0$.
Set $A_p=x^6$ and $\beta_p=x^6c_p$. Both generator images
are now their conjugates by $x^6$, so the resulting
automorphism is exactly $\operatorname{Int}(x^6)$.
It fixes $t$, and its coefficient map is an invertible
permutation action. Formula (D.39.20) is one formula
over $\mathbb Z[1/2]$; it simultaneously supplies all
odd-prime components at every precision.

The two-adic and odd-prime conjugators have the same
images $A_H,\beta_H$ in $H$. Isomorphism (D.39.9) therefore
assembles them into actual conjugators and an automorphism
of $\mathcal H$. Equations (D.39.16) ensure their equality
with the D.33 conjugators in $C_8$. The two component
automorphisms agree there on $s,r$, hence on the whole
common quotient. Their product consequently gives an
automorphism of $\mathcal T$. This proves generation
on the entire joint group and every displayed finite
reduction.

\paragraph{A single compatible gauge and the finite seed.}
The image $\mathcal F=\ker(\mathcal T\to S_3)$ has
abelianization $\widehat{\mathbb Z}^2$ in the marked
$x,y$ basis, because the source kernel in (D.39.1) is
contained in $\overline{[L,L]}\subset\widehat F'$.
For the assembled conjugator write
$\beta^{\rm ab}=(a,b)$ and set $k=b-a\in\widehat{\mathbb Z}$.
Then $\beta'=x^k\beta$ and
\[
 f=s\beta's(\beta')^{-1}\in\mathcal F'                   \tag{D.39.21}
\]
have an actual lift of $f$ in $\widehat F'$.
This follows directly by abelianizing: $s$ interchanges
$x$ and $y$, so $f^{\rm ab}=(b-a-k,a+k-b)=0$.
The normalized automorphism is
$\operatorname{Int}(x^k)\alpha$. It still fixes $t$;
its torsion relations give the two $m=0$ hexagons as
in D.25. Thus one common profinite gauge works at all
levels.

There is also an exact description of its odd components.
In the $p$-adic abelianization, (D.39.20) gives equal
coordinates for $c_p$, hence $k_p=-6$ for every odd
$p$. The recorded finite two-adic seed gives
$k_2+6\equiv0\pmod4$. In $\mathcal H_p$, the closed
subgroup generated by $x$ is an extension of its
order-twelve base image by a pro-$p$ group. Therefore
$x^{k+6}=1$: its $p$-component exponent is zero, and
the prime-to-$p$ part of the order divides the
prime-to-$p$ part of $12$, on which $k+6$ also vanishes.
It follows that the normalized odd component is
$\operatorname{Int}(x^{k+6})=1$ on the whole group.

For a literal finite check, combine the retained
two-adic coefficients modulo $16$ with (D.39.20)
modulo three by the Chinese remainder map to
$\mathbb Z/48$. The resulting $2\,676\,134$-letter cusp
word vanishes in all $6912$ coordinates. Substitution
on every one of the $1153$ kernel basis words agrees
with the computed graph matrix. Its mod-two rank is
$1153$, and its mod-three matrix is exactly the
invertible inner action by $x^6$, with its inverse
checked. Both $m=0$ hexagons are checked separately.
The integral lift above retains these mod-sixteen
coefficients, and the odd formula retains their
mod-three coefficients. Thus this entire $J_{48}$
representative remains unchanged in the limit.

Its literal $\beta$ word has abelianization
$(1482,1044)$ and gauge $-438$. Every lift with its
$J_{48}$ image changes these coordinates by elements
of $576\widehat{\mathbb Z}^2$, since the kernel is
$\overline{[L,L]L^{48}}$ and $L\subset F_{12}$.
Consequently the global $k$ reduces to $-438$ modulo
$576$. The loop $x$ has order $576$ in $J_{48}$,
as checked by its $576$th, $288$th and $192$nd powers.
It has order $48$ in $T_3$. Thus the same finite
gauge acts on the joint seed. In particular the
finite coset of $f$ contains a discrete commutator
word: the residual abelianization of a compatible
$F$ lift is a multiple of $(576,-576)$, removable
by powers of $x^{576}$ and $y^{576}$ in the joint
kernel.

The difference from D.38's gauge is
$-438-(-774)=336=7\cdot48$. The normalized $T_3$
actions therefore agree exactly, proving (D.39.6).
Projection of the two new normalized generator
images to $G_6$ also agrees exactly with D.38;
this is checked in all $97$ coefficient coordinates.
On the whole $J_3$ module each generator instead
has the identity action. Comparing with D.38 gives
$4544$ nonzero mod-three edge-coordinate differences
for each of $s,r$. Both $J_{12}$ images are replaced
as well. These are comparisons under the actual
quotient maps, including the common gauge.

The normalized $390\,284$-letter $F$ word is evaluated
on the cyclic-three, cyclic-four and cyclic-twelve
covers of D.36. All Schreier substitutions and both
compact projection maps are checked. They give
$-I$ on cubic four-torsion and $I$ on quartic
three-torsion. For the cubic Tate-sixteen refinement,
the retained $T_3$ map and the exact basis matrix
$\left(\begin{smallmatrix}-1&1\\-1&0\end{smallmatrix}\right)$
of D.38 give precisely (D.39.7). The preserved $G_6$
map gives the separately marked fixed-fibre pair.
D.36's arithmetic graph excludes the pair in
(D.39.7), and D.25 excludes the fixed-fibre pair
from the Galois image.

For every $m$ divisible by $48$, the common quotient
of $J_m$ and $T_3$ is still $C_8$, by (D.39.13) and
the relative-prime factorization. The automorphism
preserves every subgroup $m\widehat{\mathbb Z}^{1153}$
of the coefficient kernel. Hence its reductions
to $J_m\times_{C_8}T_3$ are compatible automorphisms.
Their pulled-back kernels under $B_3\to B$ are
normal, finite-index, and contained in $PB_3$;
the braid centre is killed. The common commutator
and torsion/cusp equations give the finite GT0-shadow
conditions of \cite[Definitions 2.1 and 2.3]{dBDHP2026}.
The seed cusp has order $1152$, since $x=t^{-2}$
has order $576$ and $t$ has even order. Its two
pure loops therefore have order $576$ and the
pulled-back kernel has $N_{\rm ord}=576$.

Finally $|J_{48}|=6912\,48^{1153}$ and
$|T_3|=18\,2^{11568}$. Dividing their product by
(D.39.3) proves every order in (D.39.5). The same
normal-kernel argument proves the exact source
maps in the statement. The word
$[a_1^{12},a_2^{12}]^2$ lies in $[L,L]$ and has
the nonzero degree-sixteen Magnus term computed
in D.38, so it also belongs to $D_9K$ and survives
in the full free pro-two source. The canonical
map to D.38 has kernel
\[
 \frac{\overline{[L,L]L^{12}}\cap\overline{D_9K}}
      {\overline{[L,L]}\cap\overline{D_9K}}.                \tag{D.39.22}
\]
It contains the nonzero class of $x^{144}$: that
element is killed by both finite components of
D.38, but has abelianization $(144,0)$ in the
present $\mathcal F$. Its survival is also checked
in the new $J_{48}$ seed. These exact kernels
specify the source of the lifting test; D.45 proves that
its fibre over this recorded action is empty.
\end{proof}
